\begin{problem}[Legacy AG3 Problem 11: the basic open subscheme]\label{prob:vi18-p01}
Let $A$ be a ring, let $f\in A$, and put $X=\Spec A$.  Show, as locally ringed spaces, that
\[
\bigl(D(f),\OO_X|_{D(f)}\bigr)\cong\Spec A_f.
\]
\end{problem}

\ifdefined\FullProblemDossiers
\paragraph{Why this problem matters.}
This is the exact point where localization becomes geometry.  It upgrades the formula
$\OO_X(D(f))=A_f$ to an identification of the entire open locally ringed space with a new
affine spectrum.  Distinguished opens can therefore serve as genuine affine charts.

\paragraph{Useful ideas.}
Keep the proof in four layers: points, topology, sheaves, then local rings.  Omitting any
layer proves less than an isomorphism of locally ringed spaces.

\paragraph{Guided hints.}
\begin{enumerate}
\item Use extension--contraction for the localization $A\to A_f$.
\item Track $D(g/1)\subseteq\Spec A_f$ under contraction.
\item Compare sections using $(A_f)_{g/1}\cong A_{fg}$.
\item Check compatibility with every further restriction between basic opens.
\item At $\mfp\in D(f)$ compare $(A_f)_{\mfp A_f}$ with $A_{\mfp}$ and their maximal ideals.
\end{enumerate}
\fi

\begin{solution}
We recover every layer of the legacy argument explicitly.

\emph{Step 1: prime ideals and points.}
Let $S=\{1,f,f^2,\dots\}$.  If $\mfp\in\Spec A$ and $f\notin\mfp$, then
$S\cap\mfp=\varnothing$, so $\mfp A_f$ is a prime ideal of $A_f$.  Conversely, if
$\mfq\in\Spec A_f$, then $\mfq\cap A$ is prime and cannot contain $f$, because $f/1$ is a
unit in $A_f$.  The localization correspondence therefore gives inverse bijections
\[
\Spec A_f\longrightarrow D(f),
\qquad
\mfq\longmapsto\mfq\cap A,
\]
and
\[
D(f)\longrightarrow\Spec A_f,
\qquad
\mfp\longmapsto\mfp A_f.
\]
Write the first map as $\theta$.

\emph{Step 2: the Zariski topology.}
A basic open in $\Spec A_f$ has the form $D(g/1)$ with $g\in A$.  Under $\theta$,
\[
\theta\bigl(D_{\Spec A_f}(g/1)\bigr)
=
\{\mfp\in\Spec A:f\notin\mfp,\ g\notin\mfp\}
=
D(f)\cap D(g)
=
D(fg).
\]
The sets $D(g/1)$ form a basis on $\Spec A_f$, while the sets $D(fg)$ form a basis for
the subspace topology on $D(f)$.  Thus $\theta$ is a homeomorphism.

\emph{Step 3: structure sheaves on the corresponding bases.}
Let
\[
V'=D_{\Spec A_f}(g/1)
\qquad\text{and}\qquad
V=D(fg)\subseteq D(f)
\]
be corresponding basic opens.  Then
\[
\OO_{\Spec A_f}(V')=(A_f)_{g/1},
\]
whereas
\[
(\OO_X|_{D(f)})(V)=\OO_X(D(fg))=A_{fg}.
\]
Associativity of localization gives a canonical isomorphism
\[
\Phi_g:(A_f)_{g/1}\xrightarrow{\sim}A_{fg}.
\]
On fractions it is
\[
\Phi_g\left(\frac{a/f^m}{(g/1)^n}\right)
=
\frac{a}{f^mg^n}.
\]
The inverse is forced by the universal property: the composite
$A\to(A_f)_{g/1}$ sends both $f$ and $g$ to units, and therefore factors uniquely through
$A_{fg}$.

\emph{Step 4: compatibility with restriction.}
Suppose $D(fh)\subseteq D(fg)$.  On the $A$ side the restriction map is the canonical map
\[
A_{fg}\longrightarrow A_{fh}.
\]
On the $A_f$ side it is the corresponding restriction between
$(A_f)_{g/1}$ and $(A_f)_{h/1}$.  Because all these maps are induced from $A$ by universal
properties of localization, the square formed by
\[
(A_f)_{g/1}\xrightarrow{\Phi_g}A_{fg},
\qquad
(A_f)_{h/1}\xrightarrow{\Phi_h}A_{fh}
\]
and the two restriction maps commutes.  Hence the $\Phi_g$ define an isomorphism of the two sheaves on a basis, and
therefore a unique isomorphism of sheaves on all opens.

\emph{Step 5: stalks and the local-ring condition.}
Let $\mfp\in D(f)$ and let $\mfq=\mfp A_f$.  Then
\[
\OO_{\Spec A_f,\mfq}
=
(A_f)_{\mfp A_f}
\cong
A_{\mfp}
=
\OO_{X,\mfp}.
\]
Indeed, localizing at $\mfp$ already inverts $f$ because $f\notin\mfp$.  Under the
isomorphism the unique maximal ideal on both sides is identified with
$\mfp A_{\mfp}$.  Thus the stalk map is an isomorphism of local rings, not merely of rings.

Combining the point bijection, homeomorphism, compatible sheaf isomorphism, and local stalk
isomorphisms gives
\[
\boxed{
\bigl(D(f),\OO_X|_{D(f)}\bigr)\cong\Spec A_f
}
\]
as locally ringed spaces.
\end{solution}

\ifdefined\FullProblemDossiers
\paragraph{What happened mathematically.}
The same localization operation appears simultaneously in the point correspondence, the
basis of opens, the section rings, and the stalks.  The theorem works because all four
manifestations are compatible.

\paragraph{Extensions.}
For a multiplicatively closed set $S\subseteq A$, the same argument identifies
$\Spec S^{-1}A$ with the primes of $A$ disjoint from $S$, equipped with the corresponding
restricted structure sheaf.  Principal localization is the one-element case.

\paragraph{Related problems.}
Problems~\ref{prob:vi18-p03} and \ref{prob:vi18-p04} isolate the topology and stalk
components.  The universal property for morphisms into affine schemes is AG3 Problem 12 and
remains reserved for VI/19.

\paragraph{Provenance.}
\textbf{Legacy AG3 Problem 11, recovered in full from
\texttt{theory-of-algebraic-geometry-3.tex}; reserved by the Volume VI ledger for VI/18 and
not consumed in VI/17.}
\else
\paragraph{Provenance.} \textbf{Legacy AG3 Problem 11; full canonical destination VI/18.}
\fi
